% Part: computability
% Chapter: recursive-functions
% Section: notations-pr-functions

\documentclass[../../../include/open-logic-section]{subfiles}

\begin{document}

\olfileid{cmp}{rec}{not}
\olsection{ఆదిమ పునరావృత్తి సంకేతనాలు}

ఆదిమ పునరావృత్త ప్రమేయాలకు ఖచ్చితమైన ఆగమనాత్మక వర్ణన ఉండడం వల్ల,
వాటిని క్రమబద్ధంగా వర్ణించగలం. ఉదాహరణకు, అలాంటి ప్రతి ప్రమేయానికీ ఒక
``సంకేతనం''ను కింది విధంగా కేటాయించవచ్చు. శూన్య ప్రమేయం, ఉత్తరగామి,
ప్రక్షేపాలకు వరుసగా $\Zero$, $\Succ$, $\Proj{n}{i}$ సంకేతాలను
వాడండి. ఇప్పుడు $h$ను $k$-స్థానిక ప్రమేయం~$f$, $n$-స్థానిక
ప్రమేయాలు $g_0$, \dots,~$g_{k-1}$ల నుంచి సంయుక్తం ద్వారా
నిర్వచించామనీ, ఆ తరువాతి ప్రమేయాలకు $F$, $G_0$, \dots,~$G_{k-1}$
సంకేతనాలను కేటాయించామనీ అనుకుందాం. అప్పుడు కొత్త సంకేతం
$\fn{Comp}_{k,n}$ను ఉపయోగించి $h$ ప్రమేయాన్ని
$\fn{Comp}_{k,n}[F,G_0,\dots,G_{k-1}]$తో సూచించవచ్చు.

ఆదిమ పునరావృత్తి ద్వారా నిర్వచించిన ప్రమేయాలకు సదృశ సంకేతనాలను
వాడవచ్చు. $(k+1)$-స్థానిక ప్రమేయం~$h$ను $k$-స్థానిక ప్రమేయం~$f$,
$(k+2)$-స్థానిక ప్రమేయం~$g$ల నుంచి ఆదిమ పునరావృత్తి ద్వారా
నిర్వచించామనీ, $f$,~$g$లకు కేటాయించిన సంకేతనాలు వరుసగా $F$,~$G$
అనీ అనుకుందాం. అప్పుడు $h$కు కేటాయించే సంకేతనం $\fn{Rec}_k[F,G]$.

కూడిక ప్రమేయాన్ని ఆదిమ పునరావృత్తి ద్వారా ఇలా నిర్వచించామని
గుర్తుచేసుకోండి:
\begin{align*}
  \Add(x_0, 0) & = \Proj{1}{0}(x_0) = x_0\\
  \Add(x_0, y+1) & = \Succ(\Proj{3}{2}(x_0, y, \Add(x_0, y))) = \Add(x_0, y) +1
\end{align*}
ఇక్కడ $f$ పాత్రను $\Proj{1}{0}$ పోషిస్తుంది; $g$ పాత్రను
$\Succ(\Proj{3}{2}(x_0,y,z))$ పోషిస్తుంది. $1$-స్థానిక ప్రమేయం
$\Succ$, $3$-స్థానిక ప్రమేయం $\Proj{3}{2}$ల నుంచి సంయుక్తంతో ఇది
లభిస్తుంది కనుక దీనికి $\fn{Comp}_{1,3}[\Succ,\Proj{3}{2}]$ అనే
సంకేతనం కేటాయిస్తాం. ఈ ఏర్పాటుతో కూడిక ప్రమేయాన్ని
\[
\fn{Rec}_1[\Proj{1}{0},\fn{Comp}_{1,3}[\Succ,\Proj{3}{2}]].
\]
అని సూచించవచ్చు. ఆదిమ పునరావృత్త ప్రమేయాలను లెక్కించేటప్పుడు వంటి
సందర్భాల్లో ఈ సంకేతనాలు ఉపయోగపడతాయి.

\begin{prob}
  $\Mult$కు సంపూర్ణ ఆదిమ పునరావృత్త సంకేతనాన్ని ఇవ్వండి.
\end{prob}

\end{document}
